\section*{R\'esum\'e}
\addcontentsline{toc}{section}{R\'esum\'e}

Ce rapport retrace la conception, l'impl\'ementation et l'\'evaluation d'une
recherche MCTS multi\oe{}ur dans LapZero-Chess, un moteur d'\'echecs de type
AlphaZero entra\^in\'e sur une machine personnelle. Avant ce chantier, le
moteur savait d\'ej\`a regrouper plusieurs positions dans une seule inf\'erence
GPU gr\^ace \`a un virtual loss de type Leela Chess Zero, mais un seul c\oe{}ur
CPU pr\'eparait les descentes. L'objectif \'etait de faire pr\'eparer ces
descentes par plusieurs c\oe{}urs sur un arbre partag\'e, sans corruption
silencieuse et sans perte mesurable de qualit\'e de recherche.

La solution retenue est une recherche \emph{par vagues} : des collecteurs
persistants descendent en parall\`ele avec des plateaux priv\'es, d\'eposent
leurs feuilles sous r\'eservation atomique, s'arr\^etent \`a une barri\`ere, puis
le coordinateur ex\'ecute une seule inf\'erence group\'ee, d\'eveloppe les
feuilles et effectue les remont\'ees de valeur. La correction repose sur des
\'etats de n\oe{}ud atomiques, une garde RAII qui poss\`ede le chemin r\'eserv\'e,
une table de transposition par snapshots et verrous ray\'es, et un budget
exact de simulations. La production continue, qui ferait descendre les
workers pendant les inf\'erences, est rest\'ee hors p\'erim\`etre : elle fera
l'objet d'une d\'ecision s\'epar\'ee.

Le gain mesur\'e est modeste mais reproductible : de $+7\,\%$ \`a $+26\,\%$ de
simulations par seconde selon la position, avec un plafond d\`es deux workers,
un p95 de latence qui ne se d\'egrade pas, et une qualit\'e non inf\'erieure sur
2\,500 puzzles. La raison de ce plafond est claire dans les mesures :
l'inf\'erence r\'eseau occupe 96 \`a 99\,\% du temps, donc parall\'eliser la
pr\'eparation CPU ne peut pas faire davantage. Le r\'eglage
\code{MCTS\_WORKER\_COUNT = 8} a \'et\'e activ\'e dans le bot UCI.

\vspace{0.6cm}

\begin{essentiel}
\textbf{L'essentiel en dix points.}
\begin{enumerate}[leftmargin=1.4em,itemsep=2pt]
  \item Avant le multi\oe{}ur, chaque simulation payait une latence fixe
  d'environ 3,4\,ms, et le d\'ebit \'etait plat (286 \`a 299 simulations par
  seconde) quelles que soient la position et la profondeur. Le batching par
  virtual loss a fait monter ce d\'ebit \`a 320--850 selon la position et la
  taille de lot, et un lot de 8 a \'et\'e valid\'e pour la qualit\'e.
  \item Apr\`es le batching, l'inf\'erence repr\'esente 96 \`a 99\,\% du temps
  de recherche. Le goulot restant est m\'ecanique : un seul c\oe{}ur pr\'epare
  les descentes de chaque lot, et les lots partiels co\^utent cher en latence.
  \item L'architecture livr\'ee s\'epare strictement deux phases : collecte
  parall\`ele sans \'ecriture des statistiques d'arbre, puis traitement
  s\'equentiel apr\`es barri\`ere (inf\'erence unique, expansion, remont\'ees).
  \item La propri\'et\'e d'une feuille passe par une transition atomique
  \code{Unexpanded -> Pending}. Une seule descente gagne ; les autres d\'eclarent
  une collision, rendent leur permis et recommencent.
  \item Le virtual loss devient visible pendant la s\'election, pas apr\`es la
  collecte, et une garde RAII le restaure en cas de collision, d'exception ou
  de lot partiel. Aucun chemin d'erreur ne d\'epend d'une boucle manuelle.
  \item La table de transposition ne rend plus jamais de pointeur vers une
  entr\'ee modifiable : chaque lecture copie un \emph{snapshot} sous verrou de
  bande, index\'e sur l'index physique r\'eel.
  \item Le budget est exact : chaque simulation demand\'ee produit exactement
  une remont\'ee de valeur, et aucune r\'eservation ne survit au repos.
  \item La validation combine tests d\'eterministes avec faux r\'eseau, points
  de crochet pour forcer les entrelacements, tests de mutation, stress de
  1\,800 recherches et self-play partag\'e. Sous MSVC, faute de
  ThreadSanitizer, cette discipline est le filet principal.
  \item Mesures : m\'ediane de d\'ebit $+8$ \`a $+26\,\%$ selon la position pour
  8 workers, p95 de latence sous $+5\,\%$ partout, et qualit\'e non
  inf\'erieure sur 2\,500 puzzles (delta $-0,28$ point, IC95
  $[-0,8\,;+0,24]$).
  \item La production continue, l'assainissement de l'appel \`a l'\'evaluateur
  et le cas de l'arbre chaud restent des chantiers s\'epar\'es, document\'es en
  fin de rapport.
\end{enumerate}
\end{essentiel}

\vspace{0.4cm}

\noindent\textbf{Mots-cl\'es :} MCTS, recherche parall\`ele, arbre partag\'e,
virtual loss, table de transposition, C++17, ONNX Runtime, AlphaZero,
test-driven development.

\subsection*{Comment lire ce document}

Le rapport suit un ordre narratif. Le chapitre~1 pose le probl\`eme et la
question. Le chapitre~2 d\'ecrit l'\'etat du moteur juste avant le chantier,
avec les mesures qui ont motiv\'e la conception. Le chapitre~3 est le c\oe{}ur
conceptuel : chaque m\'ecanisme y est justifi\'e, alternatives \'ecart\'ees
comprises. Le chapitre~4 raconte l'impl\'ementation en dix t\^aches, avec les
pi\`eges rencontr\'es et les extraits de code d\'ecisifs. Le chapitre~5 d\'ecrit
la strat\'egie de validation. Le chapitre~6 pr\'esente toutes les campagnes de
mesure. Le chapitre~7 discute les r\'esultats, leurs limites et les le\c{c}ons.
Le chapitre~8 conclut et ouvre les perspectives. Les annexes contiennent la
chronologie des commits, l'inventaire des fichiers, les commandes de
reproduction et un glossaire.

Les encadr\'es bleus r\'esument chaque section importante, les encadr\'es ambre
signalent des pi\`eges v\'ecus, et les encadr\'es gris d\'ecrivent les
protocoles de mesure. Les valeurs num\'eriques proviennent des fichiers
\code{out/multicore/} conserv\'es hors du d\'ep\^ot (ils sont volumineux) ; les
rapports agr\'eg\'es cit\'es, eux, sont commit\'es.
